%%
% 中英文摘要
%%

\begin{abstract}
本文面向树莓派5等无GPU ARM边缘平台上的LeRobot在线控制任务，研究 ACT 策略在 SO-101 六自由度机械臂中“观测--推理--执行”闭环的系统级性能瓶颈。随着端到端机器人学习框架逐渐进入低成本硬件平台，系统能否稳定达到目标控制频率，不仅取决于模型规模和硬件配置，也受到 Python 运行时、深度学习框架线程行为、Linux 内核调度以及设备 I/O 路径的共同影响。本文以在线控制主循环为切入点，构建覆盖推理框架、Python 运行时、Linux 内核和硬件 I/O 的分层分析方法。实验结合阶段计时、ftrace、strace 和内核插桩，对摄像头读取、舵机通信、ACT 前向推理和帧率控制等环节进行量化分析，并对测量工具自身扰动进行校准，以区分真实瓶颈与观测手段引入的额外开销。

实验结果表明，在 30 FPS 目标配置下，同步系统实测帧率约为 18.7 FPS，主要瓶颈来自 ARM CPU 上的 ACT 前向推理。工具校准实验显示，延迟导出的 ftrace 对主循环扰动较小，更适合后续定量分析；进一步的模块分解表明，单次完整推理平均耗时约 1925 ms，其中 ResNet18 视觉编码和 Transformer 模块占主要比例。相比之下，摄像头异步预取后主线程观测开销较低，舵机写入系统调用仅为微秒级，执行阶段主要受串口波特率和半双工总线约束。针对推理阻塞，本文验证了异步推理方案。当触发阈值取 0.30 时，系统有效控制频率达到 27.42 FPS，相比同步基线提升 70.0\%，并显著减少动作队列耗尽造成的停顿。结果说明，系统级分层剖析能够有效定位边缘机器人在线控制瓶颈，异步推理可在不改变模型结构的条件下显著提升控制频率，并为后续推理运行时优化、操作系统调度改进和机器人基础软件设计提供依据。
\end{abstract}

\begin{abstractEn}
This thesis studies the system-level performance bottlenecks of the LeRobot online control loop on Raspberry Pi 5, a GPU-free ARM edge platform. The target system runs an ACT policy on an SO-101 six-degree-of-freedom robotic arm and follows a perception--inference--execution loop. The online control loop is used as the entry point to build a layered profiling method covering the inference framework, Python runtime, Linux kernel, and hardware I/O. Phase timing, ftrace, strace, and custom kernel instrumentation are combined to quantify camera acquisition, servo communication, ACT forward inference, and frame-rate control.

The results show that, under a 30 FPS target configuration, the synchronous system reaches about 18.7 FPS, and the dominant bottleneck is ACT forward inference on the ARM CPU. A full inference pass takes about 1925 ms on average, with the ResNet18 visual encoder and Transformer modules accounting for most of the latency. After asynchronous camera prefetching, observation overhead on the main thread is low, while servo write system calls remain at the microsecond level. To reduce inference blocking, this thesis evaluates an asynchronous inference scheme. With the trigger threshold set to 0.30, the effective control rate reaches 27.42 FPS, improving by 70.0\% over the synchronous baseline. These results show that layered system profiling can identify the real bottlenecks of edge robotic control, and asynchronous inference can significantly improve control frequency without changing the model architecture.
\end{abstractEn}
